\section{실험 결과의 평가}

\subsection{프로토콜}

실험 스크립트는 MDLM을 자체 \texttt{ddpm\_cache} 샘플러와 비교하고, SEDD를 자체 해석적
샘플러와 비교한다. 두 실험 모두 MDLM 코드베이스의 OpenWebText 체크포인트와 시퀀스 길이 1024를
사용한다. 역방향 스텝 수는 $10$, $20$, $50$, $100$, $256$, $512$이다. 품질 지표는 다음과 같다.
\begin{itemize}
  \item GPT-2로 계산한 토큰 가중 퍼플렉서티. 낮을수록 좋다.
  \item 생성 텍스트와 참조 텍스트의 유니그램 토큰 히스토그램 사이의 Jensen--Shannon 발산.
  낮을수록 좋다. 참조 문장은 WikiText-2 원시 검증 부분집합에서 가져온 비어 있지 않은 문자열이다.
\end{itemize}

\subsection{품질 결과}

\begin{table}[t]
\centering
\small
\caption{품질 결과. $\mathrm{Gain}_{\mathrm{PPL}}
=100(\mathrm{baseline}-\mathrm{geodesic})/\mathrm{baseline}$이며, 양수이면 측지 모멘텀이
우세하다. 굵은 글씨는 각 쌍에서 더 낮은 값을 나타낸다.}
\label{tab:quality}
\begin{tabular}{llrrrrr}
\toprule
모델 & 스텝 & 기준 PPL & 측지 PPL & $\mathrm{Gain}_{\mathrm{PPL}}$ & 기준 JS & 측지 JS \\
\midrule
MDLM & 10  & 551.953 & \textbf{407.473} & $+26.18\%$ & \textbf{0.38659} & 0.38850 \\
     & 20  & 261.453 & \textbf{212.993} & $+18.53\%$ & \textbf{0.37718} & 0.38073 \\
     & 50  & 140.075 & \textbf{120.542} & $+13.95\%$ & \textbf{0.36864} & 0.37629 \\
     & 100 & 92.421  & \textbf{81.864}  & $+11.42\%$ & 0.38260 & \textbf{0.38185} \\
     & 256 & \textbf{65.431} & 65.665 & $-0.36\%$ & \textbf{0.37828} & 0.38066 \\
     & 512 & \textbf{50.905} & 51.801 & $-1.76\%$ & 0.38224 & \textbf{0.37800} \\
\midrule
SEDD & 10  & 486.276 & \textbf{434.932} & $+10.56\%$ & \textbf{0.35628} & 0.36043 \\
     & 20  & 235.867 & \textbf{214.863} & $+8.91\%$  & 0.34810 & \textbf{0.34707} \\
     & 50  & 123.169 & \textbf{114.949} & $+6.67\%$  & 0.34771 & \textbf{0.34503} \\
     & 100 & \textbf{82.088} & 84.184 & $-2.55\%$ & 0.35678 & \textbf{0.35078} \\
     & 256 & \textbf{59.954} & 63.276 & $-5.54\%$ & \textbf{0.34519} & 0.34670 \\
     & 512 & 51.695 & \textbf{50.903} & $+1.53\%$ & \textbf{0.36408} & 0.36661 \\
\bottomrule
\end{tabular}
\end{table}

\begin{figure}[t]
  \centering
  \includegraphics[width=0.94\linewidth]{quality_curves.pdf}
  \caption{퍼플렉서티 결과. 왼쪽 패널은 각 모델의 관측 범위에 맞춘 선형 세로축으로 원시
  GPT-2 퍼플렉서티를 표시하므로 적은 스텝에서의 차이를 직접 확인할 수 있다. 오른쪽 패널은
  상대 PPL 감소율 $(\mathrm{baseline}-\mathrm{geodesic})/\mathrm{baseline}$을 나타낸다.
  양의 막대는 개선, 음의 막대는 저하를 뜻한다. 두 개선률 패널은 같은 $[-8,30]$ 퍼센트포인트
  범위를 유지한다.}
  \label{fig:quality_curves}
\end{figure}

MDLM의 10, 20, 50, 100 스텝에서 측지 샘플의 GPT-2 퍼플렉서티는 각각 26.18\%,
18.53\%, 13.95\%, 11.42\% 낮다. 256 및 512 스텝에서는 차이가 사라지고 기준 방법이 약간
우세하다. 이 패턴은 적은 스텝을 보정하려는 본 방법의 의도와 일치한다.
\cref{eq:dynamic_damping}의 MDLM 감쇠는 균일 격자의
스텝 수가 50을 넘으면 $\alpha$도 감소시킨다. 그 값은 100, 256, 512 스텝에서 각각 약
$0.25$, $0.098$, $0.049$다.

SEDD는 10--50 스텝에서 더 작은 퍼플렉서티 감소를 보이고, 100--256 스텝에서는 성능이
저하되며, 512 스텝에서는 작게 개선된다. 구현은 모든 균일 격자에서 첫 스텝 이후
$\alpha=1/2$을 사용했고 MDLM의 고스텝 감쇠는 사용하지 않았다. 이것이 오버슈트에 기여했을
가능성은 있으나, 현재 실험만으로 원인을 분리할 수는 없다. Jensen--Shannon 발산은 두 모델
모두에서 결과가 혼재하며 일관된 분포 개선을 보이지 않는다.
